% Préambule
\documentclass[12pt]{ULojpv}

% Chargement des packages supplementaires
\usepackage[utf8]{inputenc}
\usepackage[T1]{fontenc}
\usepackage[french]{babel}
\usepackage{pifont}
\usepackage{comment}

% Definitions des parametres de l'en-tete
\Cours{GEL--1001 Design I (méthodologie)}             % Nom du cours
\NumeroEquipe{2}                                     % Numero de l'equipe
\NomEquipe{NordFaune}                               % Nom de l'equipe
\Objet{Procès-verbal}                                 % Nom du document
\SujetRencontre{Organisation et planification}             % Sujet de la rencontre
\DateRencontre{2026/02/12}                            % Date de la rencontre
\LocalRencontre{En ligne (Teams)}                            % Local de la rencontre
\HeureRencontre{8h30--9h20}                          % Heure de la rencontre


%--------------------------------------------------------------------------------------
%--------------------------------- corps du document ----------------------------------
%--------------------------------------------------------------------------------------
\begin{document}
\entete
\begin{enumerate}

% nouveau point
\item \textbf{Ouverture de la réunion}
Heure: 8h30

% nouveau point
\item \textbf{Nomination ou confirmation du président et du secrétaire}

\begin{tabular}{@{}ll}
   Président: Chelcie Sonon
   & Secrétaire: Louise Lesser
\end{tabular}

% nouveau point
\item \textbf{Adoption de l'ordre du jour}

L'ordre du jour proposé est adopté à l'unanimité.

% nouveau point
\item \textbf{Lecture et adoption du procès-verbal de la réunion du 5 février 2026}

Le procès verbal de la réunion du 5 février 2026 a été adopté à l'unanimité.

% nouveau point
\item \textbf{Affaires découlant du procès-verbal}
\begin{enumerate}
   \item Remise du rapport 1
   
Suite à la décision prise lors de la réunion précédente de débuter la rédaction du rapport 1, les membres ont présenté l’avancement de leur travail individuel. Nous avons notamment mis en commun les propositions d’objectifs du projet.
    

\end{enumerate}

\begin{comment}
\begin{enumerate}

\item Affaire \#1

\begin{enumerate}

\item Sous-affaire \#1

Text here

\item Sous-affaire \#2

Text here

\end{enumerate}

\item Affaire \#2

Text here

\end{enumerate}
\end{comment}

% nouveau point
\item \textbf{Points à traiter}

\begin{enumerate}

\item Retour sur le séminaire du 11 février.

Les membres de la réunion ont tous discuté sur le contenu présenté
lors du séminaire du 11 février dernier portant sur la définition du problème. Nous avons également échangé sur notre compréhension de la définition du problème et sur les éléments clés à prendre en compte pour la suite du projet.

\item Répartition des tâches

Lors de la rencontre, nous avons attribué une tâche à chaque membre de l'équipe pour la semaine à suivre. Les tâches
qui ont été attribuées se trouvent à la section 8 de ce document.

\item Remise du rapport version 1

Après la mise en commun du travail personnel sur le rapport 1, nous nous sommes mis d'accord sur les objectifs finaux que nous allons utiliser. Pour rédiger le cahier des charges plus facilement, nous avons décidé de travailler en binôme. Nous avons fait la répartition suivante :
   \begin{itemize}
   \item  Assurer l'autonomie du système : Antoine et Louis-Charles ;
   \item  Garantir la robustesse du système: Mathis et Félix-Antoine;
   \item  Assurer l'observation et la gestion de données : Phanuel et Kayanne;
   \item Réduire l'impact et le coût: Chelcie et Louise;
\end{itemize} 
Il a été décidé que le travail attribué à chaque binôme devra être réalisé pour le lundi 16 février.

\item Choix de la pondération pour les objectifs du rapport version 1

Ensemble, nous avons hiérarchiser les objectifs du projet ce qui nous a ainsi permis de déterminer la pondération des objectifs.


\end{enumerate}


% nouveau point
\item \textbf{Divers}


Nous avons décidé de lire les chapitres 3 et 4 ainsi que les commentaires du rapport fourni comme exemple pour mieux comprendre les attentes du rapport version 1.

% nouveau point
\item \textbf{Répartition des tâches}

\begin{itemize}
   \item Louise: Procès-verbal;
   \item Chelcie: Remise des livrables;
   \item Phanuel: Diagramme de gannt;
   \item Félix-Antoine: Ordre du jour;
   \item Mathis: Vérification finale;
   \item Kayanne: Vérification du français;
   \item Antoine: Direction du projet;
   \item Louis-Charles: Organisation des réunions et direction de la prochaine rencontre;
\end{itemize}

Finalement, tout le monde doit également travailler avec son binôme pour continuer la rédaction du cahier des charges et des objectifs  en vue de la remise du rapport 1 le 19 février prochain.



% nouveau point
\item \textbf{Évaluation de la réunion}

La réunion a permis à tous de connaître clairement les échéances que nous nous sommes fixées ainsi que de
s'attribuer les rôles pour la semaine à suivre. Des questions ont été posées à la fin de la rencontre et elles ont
toutes reçu des réponses satisfaisantes par un ou plusieurs membres de l'équipe.

% nouveau point
\item \textbf{Date, heure, lieu et objectif de la prochaine réunion}

\begin{tabular}{@{}lll}
   Date: 2026/02/19
   & Heure: 08h30
   &  Lieu: En ligne (Teams)
\end{tabular}
%\par


% nouveau point
\item \textbf{Fermeture de la réunion}

Heure: 09h30


% nouveau point
\item \textbf{Étaient présents}

\begin{dinglist}{"33}
   \item Antoine Robert
   \item Chelcie Laura Emmanuela Sonon
   \item Félix-Antoine Guimont
   \item Kayanne Loudiyi
   \item Louis-Charles Blais
   \item Louise Geneviève Lesser
   \item Mathis Carrier
   \item Phanuel Kouadio
\end{dinglist}

\end{enumerate}

\end{document}
